\section{Serre vanishing on every cover and polynomial high-power sections}
\label{sec:serre-h1-v055}

\status{This chapter proves the full original
\texttt{exists\_\allowbreak h1Zero\_\allowbreak tw\_\allowbreak of\_\allowbreak veryAmpleBy}:
over a Noetherian ring, every actual line bundle has vanishing first
cohomology after every sufficiently large integer twist by the supplied
very-ample bundle.  Vanishing means the original coordinate-cocycle
statement for every open cover.  No smoothness assumption is added.
The same argument also proves the original homogeneous-polynomial
description of every sufficiently high power-section, and hence both
conjuncts of \texttt{serre\_\allowbreak h1Zero\_\allowbreak and\_\allowbreak sections\_\allowbreak polynomial}.
The separate multiplication, polarized base-change, geometric rank and
square-law suppliers remain explicit.}

There are two different finite constructions.  First choose one finite
homogeneous submodule of the actual graded section module.  Its
localizations already give the actual sections on every nonempty
intersection in the prescribed cover, in every integer degree.  Next,
over a Noetherian base, take finitely many successive kernels of graded
free surjections.  Their free-module bounds have one common maximum.
These choices occur before the twist degree, before a cocycle and before
an arbitrary cover used to test vanishing.

The passage from the prescribed affine cover to every other cover is
proved separately.  An affine open has the required all-cover vanishing
by explicit denominator clearing.  Local primitives on the prescribed
affine opens differ by a cocycle on that cover.  Solving this latter
cocycle makes the local primitives agree, so they descend to the original
test cover.  Exactness on one cover is therefore connected to the literal
source predicate by an actual gluing argument.

\paragraph{Notation and inherited objects.}
Let $R$ be a commutative ring, $Y$ a scheme and
$f:Y\to\operatorname{Spec}R$.  For a line bundle $M$ and an open $V$,
write $\operatorname{Sec}(M,V)$ for its actual compatible coordinate
families.  Its $R$-action uses the original map
$c_f(V):R\to\mathcal O(V)$.  If $H$ is another line bundle, put
$M(t)=M\otimes H^t$ for every $t\in\mathbb Z$, with the actual
transition convention of ASL2.  Negative powers are dual powers.

Fix the original data $\operatorname{VeryAmpleBy}(f,H,s)$:
$s_0,\ldots,s_m$ are sections of $H$; their nonzero loci
$D_i=D(s_i)$ cover $Y$; every $D_i$ is affine; and the coordinate
ratios $r_{ij}=s_j/s_i$ give a surjection
\[
 \pi_i:R[X_0,\ldots,X_m]\longrightarrow\mathcal O(D_i),
 \qquad X_j\longmapsto r_{ij}.
                                                        \tag{SHV.1}
\]
The ratios and this evaluation use the original structure map $f$.
ASL9 supplies quasi-compactness, quasi-separatedness and affine nonempty
intersections of this cover.  ASL8, ASL12, ASL14--ASL16 supply the actual
graded module
\[
 G=\bigoplus_{t\in\mathbb Z}\operatorname{Sec}(M(t),Y),
 \qquad X_i a=s_i a,
                                                        \tag{SHV.2}
\]
its localization-to-section maps, its dehomogenizations and the finite
projective affine section modules.  FCD10 and FCD13 supply the proved
finite saturation and quotient arguments.  None of these inherited
results requires the smoothness of $Y$.

For the algebraic part let $\iota$ be any finite linearly ordered set,
$\nu=|\iota|$, $S=R[X_i:i\in\iota]$, and
$Q=\bigoplus_{t\in\mathbb Z}Q_t$ a graded $S$-module with variables
of degree one.  For $I\subseteq\iota$ put $X_I=\prod_{i\in I}X_i$
and
\[
 Q_{I,t}=(Q[X_I^{-1}])_t
 =\left\{\frac{a}{X_I^k}:\
       k\ge0,\ a\in Q_{t+k|I|}\right\}.
                                                        \tag{SHV.3}
\]
This is the source's literal \texttt{homLoc}; $Q_{\varnothing,t}=Q_t$.
Its augmented localization complex is
\[
 A^p(Q,t)=\bigoplus_{\substack{I\subseteq\iota\\|I|=p}}Q_{I,t},
 \qquad
 (d x)_J=\sum_{j\in J}(-1)^{|\{i\in J:i<j\}|}
                   x_{J\setminus\{j\}}\big|_J .
                                                        \tag{SHV.4}
\]
There are finitely many subsets, so the displayed direct sums are also
products.  The complex is zero outside degrees $0,\ldots,\nu$.
Insertion of two different indices in opposite orders gives opposite
signs, so $d^2=0$.  The ordinary complex is
$C^q(Q,t)=A^{q+1}(Q,t)$, $q\ge0$, with the same differential.
Thus ordinary degree-one exactness is augmented degree-two exactness.

\subsection{One degree bound for the whole localization complex}

% BEGIN SHV1
\paragraph{SHV1: graded localization preserves the actual short exact sequence.}
\label{claim:SHV1}
If $0\to Q'\xrightarrow{u}Q\xrightarrow{v}Q''\to0$ is a short
exact sequence of degree-preserving $S$-linear maps, then for every
$I$ and every integer $t$,
\[
 0\longrightarrow Q'_{I,t}\xrightarrow{u_I}Q_{I,t}
   \xrightarrow{v_I}Q''_{I,t}\longrightarrow0
                                                        \tag{SHV.5}
\]
is exact.  The maps commute with every localization restriction, and
therefore give a short exact sequence of the complexes in (SHV.4).
The vertex assertion holds for an arbitrary variable set and a finite
$I$; finiteness of $\iota$ is only needed for the finite sums in (SHV.4).

\begin{proof}
A homogeneous element of degree $a$ in $Q''$ has a homogeneous lift
of the same degree in $Q$: take any lift and project it to degree $a$.
The other homogeneous components map to zero, since the grading is a
direct sum.  This proves surjectivity in (SHV.5) by lifting the
numerator of a fraction.

Suppose $b/X_I^k$ maps to zero.  The equality of this fraction with
zero means that $X_I^\ell v(b)=0$ for some $\ell\ge0$.
Thus $X_I^\ell b=u(a)$ for a homogeneous $a$ of degree
$t+(k+\ell)|I|$; the homogeneous lift follows by the same projection
argument.  The fraction $a/X_I^{k+\ell}$ maps to $b/X_I^k$.
If a fraction from $Q'$ maps to zero, injectivity of $u$ applied
after multiplication by the killing power of $X_I$ proves that the
original fraction was zero.  These arguments also apply to
$I=\varnothing$, where $X_I=1$.

For $I\subseteq J$ the restriction is
$a/X_I^k\mapsto X_{J\setminus I}^ka/X_J^k$.
An $S$-linear map commutes with that formula.  Taking the finite
sums over all subsets of each cardinality proves termwise exactness
and compatibility with the signed differentials.  No tensoring by
an assumed flat coefficient algebra has been used.
\end{proof}
% END SHV1

% BEGIN SHV2
\paragraph{SHV2: the free polynomial complex and its positive-degree contraction.}
\label{claim:SHV2}
For every commutative $R$, the cohomology of $A^\bullet(S,t)$
is a finite free $R$-module in each degree.  It is zero in every
degree when $t\ge1$.

\begin{proof}
We use the Laurent-monomial calculation underlying FCD15 and spell
out its degree bound.  In the vertex $I$ a basis is
\[
 X^\alpha,\qquad
 \alpha\in\mathbb Z^\iota,\quad
 \sum_i\alpha_i=t,\quad
 N(\alpha):=\{i:\alpha_i<0\}\subseteq I .
                                                        \tag{SHV.6}
\]
To obtain such a monomial as a fraction, choose $k$ large enough that
$\alpha_i+k\ge0$ for $i\in I$.  Its numerator has those exponents
on $I$ and exponents $\alpha_i$ outside $I$, and has degree
$t+k|I|$.  Conversely every homogeneous numerator is a finite sum
of these monomials.  Their linear independence follows by clearing
one common denominator and using the usual polynomial monomial
basis.  Multiplication by $X_I$ is injective on a polynomial ring
over any $R$, because it only shifts its monomial coefficients.
This argument does not assume that $R$ is a domain.

The restrictions carry a monomial to the identical Laurent monomial.
Consequently the whole complex is a direct sum, over $\alpha$, of
complexes with basis $e_I$ for $I\supseteq N(\alpha)$ in degree $|I|$.
If $N(\alpha)\ne\iota$, choose $j_0\notin N(\alpha)$ and set
\[
 h(e_I)=
 \begin{cases}
 (-1)^{|\{i\in I:i<j_0\}|}e_{I\setminus\{j_0\}},&j_0\in I,\\
 0,&j_0\notin I.
 \end{cases}                                           \tag{SHV.7}
\]
The smaller subset still contains $N(\alpha)$.  In $dh+hd$ the
term inserting and deleting $j_0$ has coefficient one.  For
another index $j$, the two possible orders of insertion and deletion
have signs differing by one in their parity: one of $j,j_0$ lies
before the other in the order.  These two terms cancel.  Hence
$dh+hd=1$.  This includes the augmented degree-zero term.

Only vectors with $N(\alpha)=\iota$ can contribute cohomology;
each gives a single copy of $R$ in degree $\nu$ and no other term.
For $\nu>0$, write $\alpha_i=-1-\beta_i$.  The possibilities are
the nonnegative integer solutions of
$\sum_i\beta_i=-t-\nu$, a finite set for each fixed $t$.
If $\nu=0$, there is one empty exponent vector when $t=0$ and
none otherwise.  Thus the surviving direct sum is finite free.
Direct sums cause no completion issue: every cochain has finite
monomial support, and the contraction preserves that support.

For $t\ge1$, an exponent vector with every entry negative cannot
have sum $t$.  When $\iota$ is empty its sum is zero, also
impossible.  Thus all summands contract and the complex is exact
in every degree.  This proves both assertions, including the
empty variable set.
\end{proof}
% END SHV2

% BEGIN SHV3
\paragraph{SHV3: an explicit common bound for all shifts of a finite free module.}
\label{claim:SHV3}
Let $F=\bigoplus_{\ell=1}^r S e_\ell$ with $\deg e_\ell=d_\ell$
in $\mathbb Z$.  Every $H^p(A^\bullet(F,t))$ is finite over $R$.
The single bound
\[
                    b(F)=1+\sum_{\ell=1}^r|d_\ell|
                                                        \tag{SHV.8}
\]
makes $A^\bullet(F,t)$ exact in every degree whenever $t\ge b(F)$.

\begin{proof}
Localizing and taking a homogeneous piece commute with this finite
direct sum.  Its $\ell$th summand is the complex
$A^\bullet(S,t-d_\ell)$, since division by $X_I^k$ leaves the
total degree of $X^\alpha e_\ell$ equal to
$\sum_i\alpha_i+d_\ell$.  The restrictions and signs are identical.
SHV2 gives finiteness for each summand.

For every $\ell$,
$d_\ell\le|d_\ell|\le\sum_j|d_j|$.  Therefore
$t\ge b(F)$ implies $t-d_\ell\ge1$, and SHV2 contracts all the
summands.  This is one bound for every cohomological degree, not a
bound chosen after that degree.  If $r=0$, the module and complex
are zero; the displayed bound is still valid.
\end{proof}
% END SHV3

% BEGIN SHV4
\paragraph{SHV4: finitely many graded kernels can be chosen before the degree.}
\label{claim:SHV4}
Suppose $R$ is Noetherian and $Q$ is finite over $S$.
For every $a\ge1$ there are finite graded free modules
$F_0,\ldots,F_{a-1}$ and finite graded modules
$Q_0=Q,Q_1,\ldots,Q_a$, with short exact sequences
\[
 0\longrightarrow Q_{j+1}\longrightarrow F_j
       \longrightarrow Q_j\longrightarrow0
       \qquad(0\le j<a).                               \tag{SHV.9}
\]
All these objects can be fixed before any twist degree or
cohomological degree is specified.

\begin{proof}
Replace a finite generating family of $Q$ by the finitely many
homogeneous components of all its members.  They still generate,
because their span contains the original generators; they belong to
$Q$, so their span cannot be larger than $Q$.  Declare a free basis
element to have the degree of the corresponding generator.  Mapping
that basis element to the generator gives a degree-preserving
surjection $F_0\to Q$.

Its kernel is graded.  Indeed, if the image of a sum of homogeneous
components is zero, each image is zero, since the target grading
is a direct sum.  Hilbert's basis theorem makes $S$ Noetherian,
and a submodule of the finite free module $F_0$ is finite.
Its homogeneous components again give a finite homogeneous
generating family.  Repeat this construction $a$ times.
This is a finite recursion depending only on the original graded
module and the chosen finite families.  It does not assert that the
last kernel is zero or projective, and it does not assume finite
global dimension for $R$.
\end{proof}
% END SHV4

% BEGIN SHV5
\paragraph{SHV5: the exactness transfer used in dimension shifting.}
\label{claim:SHV5}
Let $0\to B^\bullet\xrightarrow{u}E^\bullet
\xrightarrow{v}D^\bullet\to0$ be a termwise short exact sequence
of cochain complexes.  If $E^\bullet$ is exact in degree $p$
and $B^\bullet$ is exact in degree $p+1$, then $D^\bullet$
is exact in degree $p$.

\begin{proof}
Let $z\in D^p$ be a cocycle.  Lift it to $e\in E^p$.
Since $v(de)=d z=0$, write $de=u(b)$ with $b\in B^{p+1}$.
The equality $u(db)=d^2e=0$ and injectivity of $u$ show that
$b$ is a cocycle.  Exactness at $p+1$ gives $b=db_0$ for
$b_0\in B^p$.  Thus $e-u(b_0)$ is a cocycle in $E^p$.
It equals $de_0$ by exactness there, and applying $v$ gives
$z=d(v(e_0))$.  This is the required coboundary.

For nonnegative complexes at $p=0$, put the degree $-1$ terms
equal to zero.  Exactness of $E$ at zero says directly that
$e-u(b_0)=0$, so $z=0$.  Thus the same argument includes the
bottom degree.  Every boundary is a cocycle because $d^2=0$,
completing the exactness assertion.
\end{proof}
% END SHV5

% BEGIN SHV6
\paragraph{SHV6: the full uniform large-degree exactness theorem.}
\label{claim:SHV6}
For every Noetherian $R$, every finite $\iota$ and every finite
graded $S$-module $Q$, there is $t_0\in\mathbb Z$ such that
\[
 t\ge t_0\quad\Longrightarrow\quad
          A^\bullet(Q,t)\text{ is exact in every degree}.
                                                        \tag{SHV.10}
\]
This is the full original
\texttt{cubeComplex\_\allowbreak grCube\_\allowbreak exactAt\_\allowbreak of\_\allowbreak finite}.

\begin{proof}
Use SHV4 with $a=\nu+1$, obtaining (SHV.9) through $Q_{\nu+1}$.
Set
$t_0=\max_{0\le j\le\nu} b(F_j)$.
For every $t\ge t_0$, SHV3 makes all the complexes
$A^\bullet(F_j,t)$ exact in every degree, and SHV1 turns
(SHV.9) into short exact sequences of complexes.

Fix any $p\ge0$.  The complex $A^\bullet(Q_{\nu+1},t)$ is
zero in degree $p+\nu+1$, because this degree exceeds $\nu$.
It is therefore exact there.  Apply SHV5 to the last short exact
sequence to obtain exactness of $A^\bullet(Q_\nu,t)$ in
degree $p+\nu$.  Applying SHV5 successively gives exactness for
$Q_{\nu-1}$ in degree $p+\nu-1$, and finally for $Q$ in degree
$p$.  Each application uses an already exact free complex.

The same finite list $F_0,\ldots,F_\nu$ and the same maximum
$t_0$ work for all $p$.  No infinite family of degree-dependent
thresholds was replaced by a maximum.  Noetherianity enters only
through the finiteness of the successive graded kernels in SHV4.
This proof also handles $\nu=0$ and arbitrary negative generator
degrees.
\end{proof}
% END SHV6

\subsection{One finite module for the actual sections in every degree}

% BEGIN SHV7
\paragraph{SHV7: finite chart saturation gives all nonempty localized vertices.}
\label{claim:SHV7}
Let $G$ be any graded $S$-module over an arbitrary commutative $R$.
Suppose $G/(X_i-1)G$ is finite over $S$ for every $i\in\iota$.
There is one finite homogeneous submodule $K\subseteq G$ such that
for every integer $t$ and every nonempty $I\subseteq\iota$ the
inclusion induces a bijection
\[
                    K_{I,t}\longrightarrow G_{I,t}.
                                                        \tag{SHV.11}
\]
These bijections are the vertices of the original inclusion map
of localization cubes.  More generally, every homogeneous submodule
satisfying (SHV.12) has these bijections, without a finiteness
assumption on either module.

\begin{proof}
FCD10 gives a finite homogeneous family $g_\ell\in G$ whose
$S$-span $K$ satisfies
\[
       \forall i,\ \forall y\in G,\ \exists k\ge0:
                         X_i^k y\in K.                 \tag{SHV.12}
\]
The submodule is graded because each homogeneous component of a
polynomial multiple of a $g_\ell$ is again such a multiple.
It is finite by its displayed generating family.  This construction
makes no choice involving $t$.

Fix $i\in I$.  If $y/X_I^a$ is a homogeneous fraction, choose
$k$ as in (SHV.12).  Then
$X_I^k y=(\prod_{j\in I\setminus\{i\}}X_j^k)X_i^k y$ belongs
to $K$.  Its homogeneous degree is $t+(a+k)|I|$, so
$(X_I^k y)/X_I^{a+k}$ is a preimage in $K_{I,t}$.
For injectivity, a fraction in $K$ that is zero in $G[X_I^{-1}]$
has a numerator killed by a power of $X_I$ in $G$.  Since $K$
is an actual submodule, that equality also holds in $K$.

The maps are induced by this same inclusion; the restriction
formula following (SHV.5) proves their compatibility.  No
isomorphism at the empty vertex is asserted here: $K_t$ can be
strictly smaller than $G_t$ in small degrees.
\end{proof}
% END SHV7

% BEGIN SHV8
\paragraph{SHV8: the full finite graded module for the actual section cube.}
\label{claim:SHV8}
For the original $f,H,s,M$ satisfying
$H.\operatorname{IsLB}$, $M.\operatorname{IsLB}$ and
$\operatorname{VeryAmpleBy}(f,H,s)$ over any commutative $R$,
there is one finite graded $S$-module $K$ and, for every integer
$t$, a cube map
\[
 \Phi_t:\bigl(I\longmapsto K_{I,t}\bigr)
       \longrightarrow
       \bigl(I\longmapsto\operatorname{Sec}(M(t),D_I)\bigr),
 \qquad D_I=\bigcap_{i\in I}D_i,
                                                        \tag{SHV.13}
\]
bijective at every nonempty vertex.  This is the full original
\texttt{exists\_\allowbreak finite\_\allowbreak graded\_\allowbreak module\_\allowbreak secCube};
$K$ is chosen before all $t$.

\begin{proof}
Use the actual graded module $G$ in (SHV.2).
ASL14 gives the original $\pi_i$-semilinear dehomogenization
$G\to P_i$, where
$P_i=\operatorname{Sec}(M(0),D_i)$, together with
\[
                  G/(X_i-1)G\simeq P_i.                 \tag{SHV.14}
\]
This quotient equality is FCD13 applied to the actual map.
ASL16 makes $P_i$ finite projective over
$B_i=\mathcal O(D_i)$.  Since $\pi_i:S\to B_i$ is surjective,
any finite $B_i$-generating family also generates $P_i$ over $S$.
Thus SHV7 applies and chooses one finite homogeneous
$K\subseteq G$.

Compose the localization cube of this inclusion with the
actual cube comparison of ASL12.  On a fraction its formula is
\[
       \Phi_{t,I}\left(\frac{a}{X_I^k}\right)
              =\left(\prod_{i\in I}s_i\right)^{-k}a|_{D_I}.
                                                        \tag{SHV.15}
\]
The factors being inverted are units on $D_I$.  ASL12 proves
this formula compatible with every original restriction and
bijective for the cube of $G$; SHV7 proves bijectivity of the
first map at every nonempty vertex.  For the empty vertex the
map is the actual inclusion $K_t\subseteq G_t$.

The scalar action, the transitions of $M(t)$ and the sections
$s_i$ in (SHV.15) have not changed.  The source allows an
abstract existential module $N$; the concrete choice $N=K$
supplies all its module, grading and decomposition structures.
No Noetherian or smoothness hypothesis was needed for this
choice or its simultaneous all-integer-degree comparison.
\end{proof}
% END SHV8

% BEGIN SHV9
\paragraph{SHV9: uniform positive-degree exactness on the prescribed affine cover.}
\label{claim:SHV9}
Now assume $R$ is Noetherian.  There is one integer $t_0$ such
that for every $t\ge t_0$ the ordinary section complex of
$M(t)$ on $D_0,\ldots,D_m$ is exact in every positive degree.
In particular its degree-one exactness is the original identity
\[
 \operatorname{im}\bigl(\operatorname{Cube.d}(\operatorname{secCube})\,1\bigr)
 =\ker\bigl(\operatorname{Cube.d}(\operatorname{secCube})\,2\bigr).
                                                        \tag{SHV.16}
\]

\begin{proof}
Fix the one module $K$ and its maps from SHV8.  Apply SHV6 to
$K$ and choose its integer bound $t_0$.  At every $t\ge t_0$,
the augmented complex $A^\bullet(K,t)$ is exact in every
degree, and hence its ordinary complex is exact in degrees
$q\ge1$.  Its terms and differentials in those degrees use
only nonempty subsets of the chart index set.

The maps (SHV.13) are bijective on all these vertices and
commute with every restriction.  Taking their finite products
gives an isomorphism of the ordinary complexes with the
actual section complex.  Transport exactness through this
isomorphism.  Ordinary degree one corresponds to subsets of
cardinality two, so its incoming and outgoing maps are exactly
the two maps printed in (SHV.16).

The assertion concerns the given sections on the given cover.
The threshold was fixed from $K$ before choosing $t$ or a
cocycle.  Upgrading it to every cover requires the next
arguments, rather than a change in the meaning of cohomology.
\end{proof}
% END SHV9

\subsection{From the chosen affine cover to every original test cover}

For a line bundle $L$ on a scheme $X$, a section cocycle on opens
$(O_a)$ is a family
$\phi_{ab}\in\operatorname{Sec}(L,O_a\cap O_b)$ satisfying
\[
                   \phi_{ab}+\phi_{bc}=\phi_{ac}
                   \quad\hbox{on }O_a\cap O_b\cap O_c.
                                                        \tag{SHV.17}
\]
Its restrictions satisfy the same identity on every smaller open,
which is exactly the quantifier in the source's
\texttt{CechOne.IsCocycle}.  A primitive is a family
$\psi_a\in\operatorname{Sec}(L,O_a)$ with
$\phi_{ab}=\psi_a-\psi_b$ on intersections.
Write $\mathcal V_L(V)$ for the assertion that every section cocycle
on every open cover of $V$, with all members contained in $V$,
has such a primitive.  Cover index sets may be arbitrary.

% BEGIN SHV10
\paragraph{SHV10: a primitive on a refinement descends to the original cover.}
\label{claim:SHV10}
Let $\phi$ be a section cocycle on $(O_a)$.  Suppose
$P_k\subseteq O_{\tau(k)}$, every $O_a$ is covered by the
$P_k$, and the restricted cocycle on $(P_k)$ has a primitive
$u_k$.  Then $\phi$ has a primitive on the original $(O_a)$.

\begin{proof}
On $O_a\cap P_k$ define
\[
                   v_{a,k}=\phi_{a,\tau(k)}+u_k.
                                                        \tag{SHV.18}
\]
All terms in this and the following identities are restricted
to the displayed common open.  On
$O_a\cap P_k\cap P_l$, the cocycle identity says
$\phi_{a,\tau(l)}-\phi_{a,\tau(k)}
 =\phi_{\tau(k),\tau(l)}$, whereas the refined primitive
says $\phi_{\tau(k),\tau(l)}=u_k-u_l$.
Hence $v_{a,k}=v_{a,l}$.  The sheaf gluing property gives
one actual section $\psi_a$ on $O_a$ with these restrictions.

On $O_a\cap O_b\cap P_k$,
\[
 \psi_a-\psi_b
   =\phi_{a,\tau(k)}-\phi_{b,\tau(k)}
   =\phi_{ab}.
\]
These opens cover $O_a\cap O_b$, so equality of sections gives
the required identity there.  No finite-cover hypothesis,
separation property of $X$, or local trivialization was needed.
In particular the argument returns a primitive on the original
cover, rather than merely on some selected refinement.
\end{proof}
% END SHV10

% BEGIN SHV11
\paragraph{SHV11: section cocycles are precisely the original coordinate predicate.}
\label{claim:SHV11}
For every actual line bundle $L$,
\[
 \mathcal V_L(X)
       \quad\Longleftrightarrow\quad
             \operatorname{FirstCohomologyVanishes}(L).
                                                        \tag{SHV.19}
\]
The right side retains every original cover, chart assignment,
transition unit and coordinate function.

\begin{proof}
Write $U_p$ for the original bundle charts and $g_{pq}$ for
their units, with coordinate convention $z_p=g_{pq}z_q$.
An original test of the right side consists of opens
$V_a\subseteq U_{c(a)}$ covering $X$ and scalar functions
$z_{ab}$ on $V_a\cap V_b$ with
\[
 z_{ab}+g_{c(a)c(b)}z_{bc}=z_{ac}
                    \quad\hbox{on triple intersections}.
                                                        \tag{SHV.20}
\]
Regard $z_{ab}$ as the $c(a)$-coordinate of a section
$\phi_{ab}$ of $L$ on $V_a\cap V_b$.  The chosen bundle
chart gives a bijection between functions and sections there.
In the $c(a)$-coordinate, the section
$\phi_{ab}+\phi_{bc}$ has the left side of (SHV.20).
Thus these sections satisfy (SHV.17).
If $\mathcal V_L(X)$ holds, let $\psi_a$ be a section primitive
and let $w_a$ be its $c(a)$-coordinate.  Taking that same
coordinate on each overlap gives
\[
             z_{ab}=w_a-g_{c(a)c(b)}w_b,
                                                        \tag{SHV.21}
\]
the literal output required by the original predicate.

Conversely, let a section cocycle be given on any cover $(O_a)$.
Refine by $P_{(a,p)}=O_a\cap U_p$.  These opens cover $X$
and have the displayed original chart assignment $(a,p)\mapsto p$.
Take the $p$-coordinate of $\phi_{ab}$ on
$P_{(a,p)}\cap P_{(b,q)}$.  The section cocycle identity becomes
the original coordinate identity with the transition $g_{pq}$.
Applying the coordinate predicate gives scalar primitives on the
refinement.  The inverse chart identifications turn them into
section primitives with exactly the refined coboundary formula.
SHV10 descends those sections to the original $(O_a)$.
This proves the converse, including infinite and repeated-index
covers.
\end{proof}
% END SHV11

% BEGIN SHV12
\paragraph{SHV12: two finite powers correct a principal-cover cocycle.}
\label{claim:SHV12}
Let $V$ be affine, $B=\mathcal O(V)$, and let
$D_1,\ldots,D_r$, $D_i=D_V(f_i)$, be a finite principal cover.
For any section cocycle $c_{ij}$ of an actual line bundle $L$
on this cover, there are integers $N\ge1$, $K\ge0$ and sections
$\rho_{ik}\in\operatorname{Sec}(L,D_i)$ such that on
$D_i\cap D_j$,
\[
 f_k^K(\rho_{ik}-\rho_{jk})=f_k^{N+K}c_{ij}
                    \qquad(1\le i,j,k\le r).
                                                        \tag{SHV.22}
\]
All occurrences of $f_k$ use its original restriction from $B$.

\begin{proof}
On the affine open $D_i$ the intersection $D_i\cap D_k$
is the principal open defined by $f_k|_{D_i}$.
ASL15 extends $c_{ik}$ after a power of this function.
There are finitely many ordered pairs, so choose one $N\ge1$
at least all their exponents.  Multiplying an already chosen
extension by the extra power gives $\rho_{ik}$ with
\[
                  \rho_{ik}|_{D_i\cap D_k}=f_k^N c_{ik}.
                                                        \tag{SHV.23}
\]
For fixed $i,j,k$, the section
$\rho_{ik}-\rho_{jk}-f_k^N c_{ij}$ on $D_i\cap D_j$
restricts to zero on the intersection with $D_k$.
Indeed (SHV.23) makes that restriction
$f_k^N(c_{ik}-c_{jk}-c_{ij})$, which vanishes by the cocycle
identity.

The open $D_i\cap D_j=D_V(f_i f_j)$ is affine.  Apply the
power-killing part of ASL15 there to the section just displayed
and the function $f_k$.  A second maximum over the finitely many
triples gives one exponent $K$ killing all these discrepancies.
Expanding the resulting identity is precisely (SHV.22).
Increasing an exponent is legitimate because multiplication by
an additional power preserves a zero section.

The first power extends sections; the second makes an identity
hold on the entire pairwise overlap.  Equality only after
restriction to $D_k$ was never cancelled across a possible
zero divisor.  The argument applies over arbitrary $B$, without
a Noetherian hypothesis.
\end{proof}
% END SHV12

% BEGIN SHV13
\paragraph{SHV13: full all-cover first-cohomology vanishing on an affine open.}
\label{claim:SHV13}
For every scheme $X$, every actual line bundle $L$ and every
affine open $V\subseteq X$, the predicate $\mathcal V_L(V)$
holds.  This is the full original
\texttt{cechHOneVanishesOnOpen\_\allowbreak of\_\allowbreak isAffineOpen}.

\begin{proof}
First solve a section cocycle on a finite principal cover
$D_i=D_V(f_i)$.  The functions $f_i$ generate the unit ideal
in $B=\mathcal O(V)$: a proper ideal containing all of them
would be contained in a maximal ideal, giving a point outside
the cover.  Put $q=N+K\ge1$, with the exponents supplied
by SHV12.  Their powers $f_i^q$ also generate the unit ideal.
For an explicit reason, write $1=\sum_i b_i f_i$ and raise
this identity to $r(q-1)+1$.  In every resulting monomial
some $f_i$ occurs at least $q$ times, so grouping terms yields
\[
                         1=\sum_i a_i f_i^q
                         \quad(a_i\in B).
                                                        \tag{SHV.24}
\]
Define
$u_i=\sum_k a_k f_k^K\rho_{ik}$ on $D_i$.
All coefficients are restricted from $B$.  Equation (SHV.22)
then gives, on $D_i\cap D_j$,
\[
 u_i-u_j
   =\sum_k a_k f_k^{N+K}c_{ij}
   =c_{ij}.
\]
Thus every cocycle on this finite principal cover is a coboundary.

Now start with an arbitrary cover $(O_a)$ of $V$ and an
arbitrary section cocycle.  Principal opens form a basis of the
affine topology, so every point of $V$ has a principal
neighborhood contained in some $O_a$.  Quasi-compactness of $V$
selects finitely many such neighborhoods $D_i$, with
$D_i\subseteq O_{\tau(i)}$.  Restrict the original cocycle to
this finite principal cover and apply the construction above.
SHV10 descends its primitive to the original cover.

If $V$ is empty, all section groups are zero and the zero
family already solves every cocycle.  Thus the empty finite
cover causes no exception.  At no point was the original
cover assumed finite, affine, or subordinate to the bundle
charts; that last issue was handled by actual sections.
\end{proof}
% END SHV13

% BEGIN SHV14
\paragraph{SHV14: local all-cover vanishing plus one cover's cocycles.}
\label{claim:SHV14}
Let $(D_i)$ cover an open $V$, with each $D_i\subseteq V$.
Suppose $\mathcal V_L(D_i)$ holds for every $i$, and suppose
every section cocycle on the particular cover $(D_i)$ has
a primitive there.  Then $\mathcal V_L(V)$ holds.  The index
sets here need not be finite.

\begin{proof}
Take a section cocycle $\phi_{ab}$ on any cover $(O_a)$ of $V$.
On each $D_i$, local vanishing supplies sections
$\sigma_{i,a}$ on $D_i\cap O_a$ such that
\[
                 \phi_{ab}=\sigma_{i,a}-\sigma_{i,b}
                 \quad\hbox{on }D_i\cap O_a\cap O_b.
                                                        \tag{SHV.25}
\]
On $D_i\cap D_j\cap O_a$ form
$\sigma_{i,a}-\sigma_{j,a}$.
For two indices $a,b$ the difference of these expressions on
their common open is
$(\sigma_{i,a}-\sigma_{i,b})
 -(\sigma_{j,a}-\sigma_{j,b})=\phi_{ab}-\phi_{ab}=0$.
They therefore glue to an actual section
$\theta_{ij}$ on $D_i\cap D_j$.

On every $D_i\cap D_j\cap D_k\cap O_a$ one has
$\theta_{ij}+\theta_{jk}=\theta_{ik}$ by cancellation of
the three $\sigma$-terms.  As these opens cover the triple
intersection, $\theta$ is a section cocycle on $(D_i)$.
By the hypothesis choose $\xi_i$ with
$\theta_{ij}=\xi_i-\xi_j$.

For fixed $a$, the sections
$\sigma_{i,a}-\xi_i$ agree on
$D_i\cap D_j\cap O_a$, by the definition of $\theta$.
Glue them on the covering opens $D_i\cap O_a$ to a section
$\psi_a$ on $O_a$.  On $D_i\cap O_a\cap O_b$,
$\psi_a-\psi_b=\sigma_{i,a}-\sigma_{i,b}=\phi_{ab}$.
Sheaf equality gives this identity on all of $O_a\cap O_b$.
These $\psi_a$ solve the original test cocycle.

The proof uses only actual restrictions and the sheaf axiom;
it does not require the pairwise intersections $D_i\cap D_j$
to be affine or quasi-compact.
\end{proof}
% END SHV14

% BEGIN SHV15
\paragraph{SHV15: the exact ordered differential and its coboundary sign.}
\label{claim:SHV15}
On a linearly ordered cover $(D_i)$, degree-one exactness of
the ordinary section complex is equivalent to the assertion
that every section cocycle on that cover has a primitive.
For a finite index set this is exactly the relationship between
the two \texttt{Cube.d} maps in (SHV.16) and the source's
\texttt{IsCoboundary}.

\begin{proof}
An ordered degree-one cochain is a family $x_{ij}$ for $i<j$,
and the first two differentials are
\[
 (\delta u)_{ij}=u_j-u_i,\qquad
 (\delta x)_{ijk}=x_{jk}-x_{ik}+x_{ij}\quad(i<j<k).
                                                        \tag{SHV.26}
\]
Restrictions to the appropriate common open are implicit.
These are the signs in (SHV.4), not a newly chosen convention.

A section cocycle $\phi$ has $\phi_{ii}=0$, since its
cocycle equation with three equal indices gives
$\phi_{ii}+\phi_{ii}=\phi_{ii}$.  It also has
$\phi_{ji}=-\phi_{ij}$, by the equation with first and third
indices equal.  Taking $x_{ij}=\phi_{ij}$ for $i<j$ gives
$\delta x=0$.  Exactness supplies $u$ with $\delta u=x$.
Set $\psi_i=-u_i$.  Then
$\phi_{ij}=\psi_i-\psi_j$ for $i<j$ by (SHV.26);
the diagonal and antisymmetry identities extend it to every
ordered pair.  The minus sign is needed because the source's
primitive uses first minus second, while $\delta u$ uses
second minus first.

Conversely, let $\delta x=0$.  Define $\phi_{ij}=x_{ij}$
for $i<j$, define $\phi_{ij}=-x_{ji}$ for $j<i$, and put
$\phi_{ii}=0$.  The expression
$\phi_{ij}+\phi_{jk}-\phi_{ik}$ on a triple intersection
changes sign under either adjacent transposition of its
indices, by antisymmetry, and is zero when two indices agree.
For increasing indices it is zero by $\delta x=0$.
Sorting the three indices therefore proves (SHV.17) for
every triple.  A cocycle primitive $\psi$ gives
$x=\delta(-\psi)$, proving exactness.  This argument is also
valid for an infinite ordered index set, where the cochains
are products and each displayed differential still uses only
finitely many terms at each component.
\end{proof}
% END SHV15

% BEGIN SHV16
\paragraph{SHV16: exactness on an affine cover implies the full coordinate predicate.}
\label{claim:SHV16}
Let $L$ be an actual line bundle on any scheme $X$.
If a linearly ordered affine open cover $(D_i)$ has its
ordinary section complex exact in degree one, then
$\operatorname{FirstCohomologyVanishes}(L)$ holds.
In particular, for every finite ordered cover this proves
the full original
\texttt{firstCohomologyVanishes\_\allowbreak of\_\allowbreak exact\_\allowbreak secCube}
with its exact \texttt{Cube.d} hypothesis.

\begin{proof}
SHV13 gives $\mathcal V_L(D_i)$ for each affine member.
SHV15 translates the given exactness into the existence of
primitives for every section cocycle on the particular
cover $(D_i)$.  SHV14 now gives $\mathcal V_L(X)$.
Finally SHV11 gives the literal coordinate predicate.

Thus for each arbitrary original chart-subordinate cover
and each of its scalar cocycles, the construction returns
coordinate functions on that same cover satisfying (SHV.21).
No assumption on the cardinality of this test cover has
entered, and no generic sheaf-cohomology comparison is left
as an unproved premise.
\end{proof}
% END SHV16

% BEGIN SHV17
\paragraph{SHV17: full original Serre first-cohomology vanishing.}
\label{claim:SHV17}
Let $R$ be any Noetherian commutative ring,
$f:Y\to\operatorname{Spec}R$ any morphism, and let
$H,M$ be the actual line bundles in an original
$\operatorname{VeryAmpleBy}(f,H,s)$ presentation.
Then there is an integer $t_0$ such that
\[
 \forall t\in\mathbb Z,\quad t\ge t_0
   \ \Longrightarrow\
       \operatorname{FirstCohomologyVanishes}(M\otimes H^t).
                                                        \tag{SHV.27}
\]
This proves \texttt{exists\_\allowbreak h1Zero\_\allowbreak tw\_\allowbreak of\_\allowbreak veryAmpleBy}
at its full original scope.

\begin{proof}
Choose the single finite graded module of SHV8 and the single
bound in SHV9.  At every $t\ge t_0$, the actual section complex
of $M(t)$ on the prescribed $D_i=D(s_i)$ is exact in degree one.
The opens $D_i$ are affine and cover $Y$, and $M(t)$ is a line
bundle for every integer $t$ because its transition functions
are the original units of $M$ times the $t$th powers of those
of $H$.  SHV16 applies to these actual data and gives (SHV.27).

The choices of the finite graded submodule, its finite list of
graded kernels and the bound precede $t$.  For each such $t$
the conclusion quantifies every cover, every assignment of an
original $M(t)$-chart to a cover member, and every coordinate
cocycle on it.  The solution is on that original cover and
uses the original transition units.  Smoothness, reducedness,
integrality, a field hypothesis and a global section
base-change theorem have not been assumed.
\end{proof}
% END SHV17

\subsection{The accompanying high-power polynomial statement}

% BEGIN SHV18
\paragraph{SHV18: a saturated finite span contains every sufficiently high piece.}
\label{claim:SHV18}
Suppose $R$ is Noetherian and $G$ is a graded $S$-module,
which need not be finite.  Let $g_1,\ldots,g_r$ be homogeneous
and let $K=\sum_\ell Sg_\ell$.  Assume
\[
 \begin{split}
 &\forall i,\ \forall y\in G,\ \exists k\ge0:\ X_i^k y\in K,\\
 &\forall y\in G,\ 
       \bigl(\forall i,\ \exists k\ge0:\ X_i^k y=0\bigr)
           \Longrightarrow y=0.
 \end{split}                                           \tag{SHV.28}
\]
Then there is $t_0\in\mathbb Z$ such that $G_t\subseteq K$
for every $t\ge t_0$.  This is the full original
\texttt{mem\_\allowbreak span\_\allowbreak of\_\allowbreak saturating\_\allowbreak of\_\allowbreak le}.

\begin{proof}
The module $K$ is finite and graded.  The first condition in
(SHV.28) is precisely the saturation used in the proof of
(SHV.11), so that proof gives bijections
$K_{I,t}\to G_{I,t}$ for all nonempty $I$.
There is no need here to assume finite dehomogenized quotients:
the finite family and its saturation have already been supplied.
The bijections commute with every restriction.

Choose a bound for $K$ by SHV6.  For $t$ above it and
$y\in G_t$, the images of $y$ in the singleton localizations
form a degree-zero cocycle of the ordinary complex of $G$:
the restrictions of the same $y$ agree on every overlap.
Lift this cocycle through the isomorphism of ordinary
localization complexes to a cocycle $z$ for $K$.
Exactness of $A^\bullet(K,t)$ in augmented degree one gives
$a\in K_t$ whose singleton images are $z$.

Hence $y-a$ maps to zero in $G[X_i^{-1}]$ for every $i$.
The fraction criterion gives some $k_i$ with
$X_i^{k_i}(y-a)=0$ for each $i$.  The second condition
in (SHV.28) forces $y=a$, proving $y\in K$.
The same bound works for every $y$ in every sufficiently
large degree.  When the variable set is empty, that second
condition already forces $G=0$, so the conclusion still holds.
\end{proof}
% END SHV18

% BEGIN SHV19
\paragraph{SHV19: the actual graded sections have no common chart-power torsion.}
\label{claim:SHV19}
Let $s_i$ be sections of an actual line bundle $H$ whose
nonzero loci cover $Y$, and let $M$ be any actual line bundle.
For the graded section module $G$ in (SHV.2),
\[
       \bigl(\forall i,\ \exists k\ge0:\ X_i^k y=0\bigr)
                    \quad\Longrightarrow\quad y=0 .
                                                        \tag{SHV.29}
\]
No affine-chart or Noetherian hypothesis is needed.

\begin{proof}
Write $y$ as its finite sum of homogeneous sections $y_t$.
Since multiplication by $X_i^k$ shifts each homogeneous
degree by $k$, the equality $X_i^k y=0$ implies
$s_i^k y_t=0$ separately for every $t$.  On $D(s_i)$
the coordinate of $s_i$ is a unit, by ASL2.  Thus
$y_t|_{D(s_i)}=0$.  The opens $D(s_i)$ cover $Y$,
so the sheaf equality of actual sections gives $y_t=0$
for every component.  Their finite sum is zero.

In this argument one cancels $s_i$ only on its nonzero
locus.  There is no assertion that multiplication by $s_i$
is injective on all global sections.
\end{proof}
% END SHV19

% BEGIN SHV20
\paragraph{SHV20: a degree-zero generator in a chart gives power saturation.}
\label{claim:SHV20}
For this claim the polynomial variable set may be arbitrary;
fix one variable $X_i$.
Let $\pi:S\to B$ be a surjective ring map with $\pi(X_i)=1$,
and let $\epsilon:G\to P$ be $\pi$-semilinear, where $G$ is
graded and $P$ is a $B$-module.  Suppose every homogeneous
element of $\ker\epsilon$ is killed by some power of $X_i$.
Let $e\in G_0$ and suppose every element of $P$ is
$b\,\epsilon(e)$ for some $b\in B$.  For every $S$-submodule
$K\subseteq G$ containing $e$ and every $y\in G$, some
$X_i^k y$ belongs to $K$.

\begin{proof}
First let $y\in G_t$ be homogeneous, allowing $t<0$.
Choose $b\in B$ with $\epsilon(y)=b\epsilon(e)$ and choose
a polynomial $F$ with $\pi(F)=b$.
Choose an integer $D\ge0$ at least $t$ and the total degree
of every monomial of $F$.  This is a finite choice even
when the variable set is infinite.  Homogenize with $X_i$:
\[
 F_D=\sum_\alpha c_\alpha X_i^{D-|\alpha|}X^\alpha .
                                                        \tag{SHV.30}
\]
It is homogeneous of degree $D$ and $\pi(F_D)=\pi(F)$
because $\pi(X_i)=1$.
Both $X_i^{D-t}y$ and $F_D e$ have degree $D$ and have
the same image under $\epsilon$.
By the homogeneous-kernel hypothesis, for some $a\ge0$,
\[
 X_i^a(X_i^{D-t}y-F_D e)=0.
\]
Thus $X_i^{a+D-t}y=X_i^aF_D e$ belongs to $K$.
All exponents used as powers are nonnegative.

For a general $y$, apply this argument to its finitely many
homogeneous components.  Increase their exponents to a common
maximum, using that $K$ is an $S$-submodule, and add.
This proves the assertion without assuming $\epsilon$
surjective, $G$ finite or the base ring Noetherian.
It is the full original
\texttt{exists\_\allowbreak X\_\allowbreak pow\_\allowbreak smul\_\allowbreak mem\_\allowbreak span\_\allowbreak of\_\allowbreak dehomogenisation},
including its arbitrary containing submodule.
\end{proof}
% END SHV20

% BEGIN SHV21
\paragraph{SHV21: every high-power section is an actual homogeneous polynomial.}
\label{claim:SHV21}
For every original $\operatorname{VeryAmpleBy}(f,H,s)$ over
a Noetherian $R$, there is a natural number $T$ such that
for every natural number $t\ge T$ and every actual section
$\sigma$ of $H^t$, there is $F\in R[X_0,\ldots,X_m]$,
homogeneous of degree $t$, with
\[
 \sigma_k=
  \operatorname{eval}\bigl(c_f(U_k),(s_j)_k\bigr)(F)
                 \qquad\text{for every original }H\text{-chart }U_k.
                                                        \tag{SHV.31}
\]
This proves the full original \texttt{exists\_\allowbreak form\_\allowbreak eq\_\allowbreak sec\_\allowbreak pw}.
Together with SHV17 it proves both original conjuncts of
\texttt{serre\_\allowbreak h1Zero\_\allowbreak and\_\allowbreak sections\_\allowbreak polynomial}.

\begin{proof}
Use the actual trivial line bundle for $M$ in (SHV.2).
Let $e\in G_0$ be the section whose coordinate is $1$.
The canonical zero-twist identification identifies
$\operatorname{Sec}(M(0),D_i)$ with $B_i=\mathcal O(D_i)$,
and the degree-zero restriction formula of ASL14 sends
$e$ to $1$ on each chart.  Every element of that chart
module is therefore a multiple of $\epsilon_i(e)$.
The original evaluation $\pi_i$ is surjective and has
$\pi_i(X_i)=1$, and ASL14 supplies its homogeneous
power-kernel property.

Apply SHV20 to each chart with the single submodule $K=S e$.
It gives the first condition in (SHV.28) for every element
of $G$, while SHV19 gives the second.  The homogeneous
family consists of this one degree-zero element $e$.
SHV18 now supplies one integer $t_0$ with $G_t\subseteq S e$
for all $t\ge t_0$.  Choose a natural number $T$ with
$T\ge\max(0,t_0)$.

An actual section $\sigma$ of $H^t$ defines an element
$y\in G_t$ by retaining its coordinates in the factor
$\mathcal O_Y\otimes H^t$.  This uses the trivial bundle's
single chart and coordinate $1$, so it introduces no choice
of trivialization or of a new section.
For $t\ge T$, write $y=F'e$.
The map $S\to G$, $F'\mapsto F'e$, preserves degree.
Taking its degree-$t$ component shows
$y=F e$ for the degree-$t$ homogeneous part $F$ of $F'$.
This remains valid for $F=0$.

Finally the definition of the polynomial action gives its
coordinate formula.  Multiplication by $X_j$ multiplies by
the original coordinate $(s_j)_k$, and a coefficient $r\in R$
acts by $c_f(U_k)(r)$.  Iterating for a monomial and adding
its finitely many coefficients shows that the $k$-coordinate
of $F e$ is exactly the right side of (SHV.31), because
$e$ has coordinate one.  The coordinate of $y$ is $\sigma_k$.
Thus (SHV.31) holds on every original chart, with one
homogeneous polynomial and the original scalar map.

The bound $T$ was fixed before $t$ and before $\sigma$.
Unlike SHV17's bound, which can depend on the arbitrary
bundle $M$, this bound uses only the fixed trivial-bundle
graded construction for $f,H,s$.
SHV17 and this assertion have exactly the two quantifier
orders in the original conjunction, so pairing them gives
that conjunction without any additional geometric premise.
\end{proof}
% END SHV21

\subsection{Return to the remaining section-law arguments}

SHV17 supplies the full input named \texttt{GSL\_SERRE} in
the earlier contracts.  That input was explicitly the first,
all-cover $H^1$ conjunct; SHV21 additionally supplies the
high-power polynomial conjunct of the original Serre theorem.
Neither conclusion assumes the smooth-projective \texttt{Perf}
record.  ASL21 supplies that record independently at its stated
smoothness hypothesis; this chapter uses the earlier actual
localization and affine-section arguments before smoothness enters.

In GSL17 the Serre input can therefore be removed from the retained
premises.  Its positive-power vanishing conclusion still requires the
original polarized base-change framework, the multiplication-pullback
generated line-bundle comparison, finite/flat/surjective and etale
multiplication geometry, and descent of every invariant function on
every relevant open.  The prime and residue-class argument in GSL17
then applies exactly as proved.  Those remaining producers have not
been supplied by a theorem about high twists.

Likewise the principal rank-one theorem, geometric integrality,
torsion-indexed section decomposition, the required two-sided section
growth, and the square/translation generators remain distinct inputs
to the later global section laws.  A polynomial description in high
degree alone does not identify those dimensions or produce those
geometric comparisons.  No whole selected RH or Dirichlet-GRH
obligation is declared closed by this chapter.

\paragraph{Pinned source map.}
All source references are to
\texttt{Coco\allowbreak{}Sgt/\allowbreak{}riemann-hypothesis-lean}, commit
\texttt{7cbf7a0c32fca1badf05931879a37c3f8a2b03b4}.
Write $C$ for
\texttt{src/\allowbreak{}Development/\allowbreak{}Abelian\allowbreak{}Variety\allowbreak{}Theory/\allowbreak{}Proofs/\allowbreak{}Coherent\allowbreak{}Cohomology\allowbreak{}Proofs/\allowbreak{}},
$A$ for the sibling
\texttt{Proofs/\allowbreak{}Abelian\allowbreak{}Scheme\allowbreak{}Law\allowbreak{}Proofs/\allowbreak{}}, and $D$ for the sibling
\texttt{Definitions/\allowbreak{}Quasicoherent\allowbreak{}Cohomology/\allowbreak{}}.
The accompanying provenance gives exact source ranges and hashes.
The principal original declarations are as follows.

\begin{itemize}
\item SHV1--SHV6: $D$ \texttt{Exists\allowbreak{}Away\allowbreak{}Frac\allowbreak{}Of\allowbreak{}Mem\allowbreak{}Hom\allowbreak{}Loc.lean},\par
$D$ \texttt{Is\allowbreak{}Graded\allowbreak{}Map\allowbreak{}Linear\allowbreak{}Combination.lean}, and
$C$ \texttt{Cube\allowbreak{}Complex\allowbreak{}Graded\allowbreak{}Module\allowbreak{}Exact\allowbreak{}At\allowbreak{}Of\allowbreak{}Large\allowbreak{}Degree.lean}.
The free calculation is in
$C$ \texttt{Exists\allowbreak{}Monomial\allowbreak{}Basis\allowbreak{}Cech\allowbreak{}Cube\allowbreak{}Polynomial\allowbreak{}Grading.lean},
\texttt{Cube\allowbreak{}Complex\allowbreak{}Exact\allowbreak{}At\allowbreak{}Of\allowbreak{}Monomial\allowbreak{}Basis.lean},
\texttt{Cube\allowbreak{}Complex\allowbreak{}Poly\allowbreak{}Grading\allowbreak{}Finite\allowbreak{}Exact.lean} and
\texttt{Cube\allowbreak{}Complex\allowbreak{}Gr\allowbreak{}Free\allowbreak{}Finite\allowbreak{}And\allowbreak{}Exact\allowbreak{}At.lean}.
\item SHV7--SHV9: $A$
\texttt{Exists\allowbreak{}Finite\allowbreak{}Graded\allowbreak{}Module\allowbreak{}Line\allowbreak{}Bundle\allowbreak{}Sections.lean},
especially \texttt{exists\_\allowbreak{}finite\_\allowbreak{}gr\allowbreak{}Cube\_\allowbreak{}hom\_\allowbreak{}bijective} and
\texttt{exists\_\allowbreak{}finite\_\allowbreak{}graded\_\allowbreak{}module\_\allowbreak{}sec\allowbreak{}Cube}.
\item SHV10--SHV16: $A$
\texttt{Cech\allowbreak{}HOne\allowbreak{}Vanishes\allowbreak{}On\allowbreak{}Open\allowbreak{}Of\allowbreak{}Is\allowbreak{}Affine\allowbreak{}Open.lean},\par
from \texttt{is\allowbreak{}Coboundary\_\allowbreak{}of\_\allowbreak{}refine} through
\texttt{h1\allowbreak{}Zero\_\allowbreak{}of\_\allowbreak{}exact\allowbreak{}At\_\allowbreak{}one}; and $C$
\texttt{First\allowbreak{}Cohomology\allowbreak{}Vanishes\allowbreak{}Of\allowbreak{}Exact\allowbreak{}Sec\allowbreak{}Cube.lean}.
The literal coordinate predicate is in
\texttt{Definitions/\allowbreak{}Abelian\allowbreak{}Scheme\allowbreak{}Laws/\allowbreak{}First\allowbreak{}Cohomology\allowbreak{}Vanishes\allowbreak{}Spec\allowbreak{}Of\allowbreak{}Is\allowbreak{}Local\allowbreak{}Ring.lean}.
\item SHV17: $C$ \texttt{Serre\allowbreak{}Vanishing\allowbreak{}H1\allowbreak{}Twist.lean},
\texttt{exists\_\allowbreak{}h1\allowbreak{}Zero\_\allowbreak{}tw\_\allowbreak{}of\_\allowbreak{}very\allowbreak{}Ample\allowbreak{}By}.
\item SHV18--SHV21: $C$ \texttt{Mem\allowbreak{}Span\allowbreak{}Of\allowbreak{}Saturating\allowbreak{}Of\allowbreak{}Le.lean}
and \texttt{Exists\allowbreak{}Homogeneous\allowbreak{}Form\allowbreak{}Eq\allowbreak{}High\allowbreak{}Power\allowbreak{}Section.lean};
the final two-conjunct interface is $A$
\texttt{Serre\allowbreak{}Vanishing\allowbreak{}Very\allowbreak{}Ample.lean}.
\end{itemize}
